\section*{导言}

这期对谈表面上是在聊赵晨阳：清华本科，UCLA 读博十五个月，参与 SGLang，从开源项目走向推理引擎创业；视频发布时，RadixArk 刚官宣一亿美元种子轮融资。若只按新闻点读，很容易把它理解成“硅谷退学潮里的年轻人故事”。但整场真正值得保存的部分，不是退学本身，而是退学、SGLang、AI Coding 和开源社区共同指向的一个更深问题：当 AI 把执行成本迅速压低之后，人类还剩哪些职责必须亲自承担？

赵晨阳反复强调“做事要本质一点”。这句话贯穿全场。他谈学术圈时，批评的是拿着锤子找钉子、虚空 formulate 一个 evaluation 的惯性；谈退学时，关心的是 applied research 到底需不需要真正接受市场和工程验证；谈 AI Coding 时，强调的是代码生产变便宜后，人类更需要负责 design、review、say no 和承担后果；谈开源社区时，真正关心的也不是“贡献数量”，而是一个项目如何在 AI 生成代码泛滥后仍然维持责任边界、审美和协作秩序。

因此，这份笔记不按时间顺序复述，而按六个主题重组：推理引擎的生态位置、退学潮背后的研究范式变化、AI 时代的 research taste、AI Coding 对代码治理的冲击、AI 作为放大器带来的认知风险，以及开源社区如何在生成式代码时代重新定义贡献。

\begin{importantbox}{阅读这份笔记的三条主线}
第一条主线是“本质性”：applied research 要有 application，fundamental research 要有自洽的问题体系，不能只在虚构靶子上优化。第二条主线是“责任边界”：AI 让执行更便宜，但也让不动脑子的 PR、代码和内容更容易泛滥。第三条主线是“社区与组织”：开源项目不只是代码仓库，而是要记录贡献、筛选责任、传递信息并组织人一起做事。
\end{importantbox}

原视频没有官方字幕，本整理基于本地 Whisper 转写和逐段校读完成。正文不做逐字稿，而是尽量提炼访谈中的判断链条和可迁移观点。

